Ce travail pratique s'incrit dans l'étude de l'optique géométrique 
et plus particulièrement des instruments modélisés dans cette théorie. 
Ses objectifs sont multiples : 
\begin{enumerate}
    \item Mise en évidence et caractérisation du pouvoir séparateur de l'oeil.
    \item Etude des grandes caractéristiques optiques du microscope.
    \item Utilisation du microscope de manière concrète dans un cas appliquée à la biologie.
\end{enumerate}
La problématique que ce travail soulève est la suivante : 
Dans quelle mesure l'utilisation d'un système optique composé comme le microscope 
permet-elle de s'affranchir de la limite de résolution naturelle de l'œil humain 
pour caractériser quantitativement des objets micrométriques ?